\section{Datasheet}\label{app:datasheet}
Following the format of Gebru et al.~\cite{gebru2021datasheets}. The dataset described here is the derived collection released with this paper, not ChEMBL itself.

\paragraph{Motivation.} The collection was assembled to estimate how much of the error of a potency model is irreducible label noise, and to compare models against that floor rather than against each other alone. It exists because published benchmarks report RMSE and $R^2$ without any reference point for the noise of the labels, so a model that has reached the limit of the data is indistinguishable from one that has not. It was assembled by the author at Innopolis University; no external funding supported it.

\paragraph{Composition.} Each instance is a molecule--target--endpoint record. The release covers 41 tasks (a task is one of 40 human single-protein targets paired with one endpoint, $K_i$ or $\mathrm{IC}_{50}$; one target yields both): 30 $\mathrm{IC}_{50}$ tasks and 11 $K_i$ tasks, 210{,}519 molecule--task records over 165{,}439 distinct structures, built from 266{,}412 document-level measurements drawn from 8{,}543 distinct ChEMBL documents published between 1978 and 2025. Tasks hold 1{,}000--9{,}567 molecules (median 4{,}778). Two files per task are released: \texttt{*.molecules.csv} (molecule identifier, canonical SMILES, label in pChEMBL units, number of source documents, year of first report, earliest source document) and \texttt{*.measurements.csv} (one row per molecule and document, with the document year). The label is the mean over a molecule's documents; 1.8--18.4\% of molecules per task (median 3.6\%) have more than one, and these replicates are what the floors are estimated from. The collection is not a sample: it is every ChEMBL record passing the filters of Section~\ref{sec:data}, so it is complete with respect to those filters and not with respect to the medicinal chemistry literature. It is not self-contained --- SMILES and labels are included, but assay descriptions, protein sequences and the documents themselves are not, and must be fetched from ChEMBL by the identifiers given. Labels are noisy by construction; quantifying that noise is the point of the paper. Censored measurements (relations `$<$', `$>$') were dropped, so the potency range is truncated at both ends, and molecules appearing in 11 MoleculeNet or 40 TDC datasets were removed (Section~\ref{sec:data}), so the collection is deliberately disjoint from those benchmarks.

\paragraph{Collection process.} Records were retrieved from the public ChEMBL REST API on 10 October 2026, which reported release ChEMBL~37 of 1 May 2026, by \texttt{scripts/01\_extract\_chembl.py}. Nothing was measured or observed directly; every value is a curated transcription by the ChEMBL team of a value published in the cited document. The 59 queried targets were chosen in two passes: 29 by hand for the first version of this work, 30 by activity count over all human single-protein targets (\texttt{scripts/13\_discover\_targets.py}). Data collection spans the publication window of the source documents, 1978--2025, which is decades earlier than the collection itself. Document metadata (authors, journal, year, type) was retrieved separately for the provenance audit of Section~\ref{sec:provenance}; see \emph{personal data} below.

\paragraph{Preprocessing and labelling.} Structures were standardised with RDKit (largest fragment, neutralisation, canonical SMILES) and keyed by the connectivity block of the InChIKey, which merges stereoisomers; a stereo-aware sensitivity analysis is reported in Section~\ref{sec:ceiling}. Values for one molecule within one document are averaged, so every remaining replicate is between documents. Molecules matching the decontamination registry by InChIKey, by connectivity block, or by ECFP4 Tanimoto similarity above 0.8 were removed: 9{,}689 of 220{,}208 (4.4\%). The raw extraction is kept alongside the processed files in the repository, so every step can be re-run or replaced. All preprocessing code is in \texttt{scripts/01\_extract\_chembl.py} through \texttt{scripts/03\_ceiling.py} and \texttt{scripts/build\_registry.py}.

\paragraph{Uses.} The collection has been used for the noise-ceiling estimates, the model comparison and the document-support analysis of this paper. It is suitable for benchmarking potency regression under the splits provided, for studying label noise and provenance, and for methodological work on evaluation. Two misuses are worth naming. The floors are between-document disagreements, not instrument error, so they should not be read as the precision of any single assay, and a model reaching the floor has not been shown to be experimentally useful: nothing here substitutes for prospective validation. And results under the random split of this collection are inflated by document support (Section~\ref{sec:support}); the document-grouped split is the one to compare against. The assay confidence threshold admits orthologous and mutant constructs, which are not separated, so the collection should not be used to study species or mutant selectivity without further filtering.

\paragraph{Distribution.} {\sloppy The collection is public at \url{https://github.com/Ruskinkot1/noise-ceiling-mol} (tag \texttt{preprint-v2}), with the ChEMBL~37 snapshot archived at \url{https://doi.org/10.5281/zenodo.23572957}. ChEMBL data are distributed under the Creative Commons Attribution-ShareAlike 3.0 Unported licence; the derived data here are released under the same licence, and the code under the MIT licence. No third party holds rights in the derived files beyond those terms, and no export controls or fees apply.\par}

\paragraph{Maintenance.} The collection is maintained by the author (\texttt{r.lukin@innopolis.university}); issues and corrections go through the GitHub repository. The release is a frozen snapshot and will not be silently updated: errata and any later version will be new tags, with the existing tags left in place so that published numbers remain reproducible. Because \texttt{scripts/01\_extract\_chembl.py} queries the live API, re-running the extraction yields a later ChEMBL release and different numbers; the Zenodo archive is the reference for this paper. Others may extend the collection by adding targets to \texttt{configs/endpoints\_extension.yaml} and re-running the pipeline.

\paragraph{Personal data.} The released per-task files contain no personal data: they hold structures, numeric labels and ChEMBL document identifiers. The provenance audit of Section~\ref{sec:provenance} did use author lists, retrieved from the ChEMBL API --- published bibliographic metadata of scientific documents, not data about private individuals. Author names are not part of the release: the published provenance files (\texttt{results/}\allowbreak\texttt{full\_clean/} and \texttt{results/ext/prov/}) record only derived indicators per document pair --- whether an author, two authors or the last author is shared, and whether the journal is the same. The author strings live only in a local metadata cache that \texttt{scripts/11\_provenance\_audit.py} rebuilds from the ChEMBL API and that is not committed. Nothing in the collection identifies subpopulations, and no human subjects are involved.
